\uuid{7eY6}
\titre{Primitive}
\auteur{bodin ; croizat ; sacre}
\organisation{exo7}
\licence{CC BY-NC-SA 4.0}
\datecreate{2022-03-08}
\niveau{L1}
\module{Analyse}
\chapitre{Calcul d'intégrales}
\sousChapitre{Primitives diverses}
\difficulte{3}
\theme{Primitive}
\exo7qcmid{hilisit.equations-differentielles.009}

\contenu{
\begin{question}
\qclassification{127.11}
Parmi les fonctions suivantes, laquelle est une primitive de $\sqrt x$ sur l'intervalle $]0,+\infty[$ ?
\begin{answers}
  \bad{$2x\sqrt x$}
  \bad{$\frac 1{2\sqrt x}$}
  \bad{$x^2\sqrt x$}
  \good{$\frac 23 x\sqrt x$}
\end{answers}
\begin{explanations} La dérivée de $x\sqrt x$ est $\sqrt x+x\times \frac 1{2\sqrt x}=\sqrt x+ \frac {(\sqrt x)^2}{2\sqrt x}=\sqrt x+\frac 12\sqrt x= \frac 32\sqrt x$ donc $\frac 23x\sqrt x$ est une primitive de  $\sqrt x$. Remarquez d'ailleurs que $x \sqrt{x}$ peut aussi s'écrire $x^{3/2}$, ce qui permet d'obtenir différemment sa dérivée : $(x^{3/2})' = \frac{3}{2} x^{3/2 - 1}= \frac{3}{2} x^{1/2} = \frac{3}{2} \sqrt{x}$. Par ailleurs, les dérivées de $x^2\sqrt x$ et de $\frac 1{2\sqrt x}$ donnent respectivement $\frac{5}{2} x \sqrt{x}$ et $\frac{-1}{4x \sqrt{x}}$, qui sont donc bien distinctes de $\sqrt{x}$.
\end{explanations}
\end{question}
}
